\documentclass{article}
\usepackage{graphicx} % Required for inserting images
\usepackage{amsmath}
\usepackage[most]{tcolorbox}

\newtcolorbox{exercisebox}[1][]{
  colback=blue!3,
  colframe=blue!45!black,
  title=Exercise,
  fonttitle=\bfseries,
  arc=2mm,
  boxrule=0.8pt,
  left=2mm,
  right=2mm,
  top=1mm,
  bottom=1mm,
  #1
}

\newtcolorbox{examplebox}[1][]{
  colback=green!3,
  colframe=green!40!black,
  title=Example,
  fonttitle=\bfseries,
  arc=2mm,
  #1
}

\newtcolorbox{notebox}[1][]{
  colback=yellow!8,
  colframe=yellow!45!black,
  title=Note,
  fonttitle=\bfseries,
  arc=2mm,
  #1
}

\newcommand{\answer}[1]{%
  \medskip
  \noindent\textbf{Answer:} #1%
}

\newcommand{\w}{\w}
\newcommand{\z}{\z}

\title{Fulton-Schubert Lecture Notes}
\author{Pedro Angulo-Umana}
\date{April 2026}

\begin{document}

\maketitle

\section{Introduction}

Begin with the 3D primitive equations in log-p coordinates:
\begin{subequations}
\label{eq:linear_system}
\begin{align}
\frac{\partial u}{\partial t} - f v
    &= -\frac{\partial \Phi}{\partial x},
    \label{eq:linear_system_a} \\[0.5em]
\frac{\partial v}{\partial t} + f u
    &= -\frac{\partial \Phi}{\partial y},
    \label{eq:linear_system_b} \\[0.5em]
\frac{\partial u}{\partial x}
+ \frac{\partial v}{\partial y}
+ \frac{1}{\rho_B}
  \frac{\partial}{\partial \z}
  \left( \rho_B \w \right)
    &= 0,
    \label{eq:linear_system_c} \\[0.5em]
\frac{\partial \Phi}{\partial \z}
    &= \frac{R T}{H_s},
    \label{eq:linear_system_d} \\[0.5em]
\frac{\partial}{\partial t}
\left(
    \frac{\partial \Phi}{\partial \z}
\right)
+ N^2 \w
    &= 0.
    \label{eq:linear_system_e}
\end{align}
\end{subequations}

\begin{exercisebox}
The first two equations of Eq.~\ref{eq:linear_system} are the zonal and meridional momentum equations. What are the corresponding ``names'' for the other three?

\answer{The continuity equation, the hydrostatic balance equation, and the thermodynamic equation.}
\end{exercisebox}
We want to collapse all the vertical structure information into a single term. The vertical coordinate shows in in Eqs. \ref{eq:linear_system_c}-\ref{eq:linear_system_e}. These equations involve $u$, $v$, $w$, $T$, and $\phi$. The horizontal equations are just in terms of $u$, $v$, and $\phi$, so we should try to eliminate $w$ and $T$ by writing them in terms of the other three variables. 

\par From the thermodynamic equation we have
\begin{equation}
    U oh 
\end{equation}
\end{document}
